\documentclass{article}
\usepackage[utf8]{inputenc}

% two packages needed for the nutrition label
\usepackage{xparse}
\usepackage{booktabs}
% this allows for the easy resize box
\usepackage{graphicx}
% for links
\usepackage{hyperref}

\hypersetup{
    colorlinks=true,
    urlcolor=blue
}

% These commands create shortcuts for the nutrition label
% Modified from: https://tex.stackexchange.com/questions/118600/how-can-i-create-a-nutrition-facts-label
\newlength{\NFwidth}
\setlength{\NFwidth}{4.5in}

\NewDocumentCommand{\NFelement}{mmm}{\normalsize\textbf{#1} #2\hfill #3}
\NewDocumentCommand{\NFline}{O{l}m}{\footnotesize\makebox[\NFwidth][#1]{#2}}

\NewDocumentCommand{\NFentry}{sm}{%
  \makebox[.5\NFwidth][l]{\normalsize
    \IfBooleanT{#1}{\makebox[0pt][r]{\textbullet\ }}%
    #2}\ignorespaces}
\NewDocumentCommand{\NFtext}{+m}
 {\parbox{\NFwidth}{\raggedright#1}}

\newcommand{\NFtitle}{\multicolumn{1}{c}{\huge\bfseries Dataset Facts}}

\newcommand{\NFRULE}{\midrule[5pt]}
\newcommand{\NFRule}{\midrule[2pt]}
\newcommand{\NFrule}{\midrule}




\title{Data Card Nutrition Label Template}
\author{Jack Bandy}

\begin{document}




\begin{figure}
    \centering
    % easily resize however you like!
    \resizebox{11cm}{!}{
    
    % For organization, I like to keep the actual card in a separate file:
    \input{Card}
    
    }
    
    \caption{A data card styled after nutrition labels. Created for the \href{https://arxiv.org/abs/2105.05241}{retrospective BookCorpus datasheet}, with help from \href{https://tex.stackexchange.com/questions/118600/how-can-i-create-a-nutrition-facts-label}{Stack Exchange}. See \href{https://arxiv.org/abs/1803.09010}{Gebru et al.} for guidance on creating complete datasheets for datasets.}
    \label{fig:data_card}
\end{figure}




\end{document}
